\chapter{结论与建议}

\section{结论}

\begin{enumerate}
\item 建设项目涉及改建项目，涉及射线装置（1 台医用电子直线加速器、CT模拟定位机，使用 1枚II类射线装置 X射线）的使用。根据《放射诊疗管理规定》，建设项目建成后开展的诊疗项目为放射治疗；根据《放射诊疗建设项目卫生审查管理规定》，按照可能产生的放射性危害程度与诊疗风险，建设项目属于危害严重类放射诊疗建设项目。
\item 本项目放射工作场所拟按功能布置放疗科，在满足诊疗流程的同时又尽可能集中设置，平面布局整体布局合理，监督区和控制区划分明确，符合相关标准要求。
\item 1 间医用电子直线加速器、CT模拟定位机房工作场所各屏蔽体防护设计厚度均不低于理论计算值；根据理论计算结果，本项目放射工作人员可能受到的年有效剂量均低于国家标准要求限值及建设单位管理目标值；公众可能受到的年有效剂量低于国家标准要求限值及建设单位提出的管理目标值。
\item 本项目拟设置的防护安全设施及措施符合放射卫生相关标准要求。
\item 建设单位拟对本项目建成后的辐射源监测、放射工作场所监测和个人剂量监测以委托检测和自主监测相结合的形式进行管理，计划配备相应质控设备，符合要求。
\item 建设单位已制定应急管理组织，明确了相关职责及应急准备与响应，符合相关要求。
\item 建设单位已成立放射防护管理组织，明确了相关职责，符合相关要求；建设单位已制定了一系列放射防护相关制度，计划根据建设项目情况对现有的相关制度进行补充完善。
\item 建设项目拟配备的人员结构符合相关要求，按相关要求对所有放射工作人员进行放射防护培训、个人剂量、职业健康和档案进行管理，符合相关要求。
\end{enumerate}

综上所述，建设项目在场所布局、屏蔽防护、放射防护措施等方面进行了考虑，在后续工作过程中，严格按照本报告中相应章节描述的各方面内容进行设计和建设，并进一步落实 9.2 节提出的建议后，建设项目在放射防护方面基本可行。

\section{建议}

\begin{enumerate}
\item 落实放射防护措施的设置；落实质控设备的配置，相应设备在使用前送至有资质的机构进行检定或校准；在设备就位后医院应及时组织相关工作人员参加培训，熟练掌握相关技能，并在日常监测过程中做好相关记录。
\item 建设项目在投入使用前，建设单位应按照国家最新法律法规及标准规范的要求及本报告建议，完善相关应急处理预案、操作规程、质量保证方案等放射防护管理制度，将相关制度如操作规程、人员岗位职责、放射治疗患者须知、放射事故应急处理措施等张贴上墙。
\item 新增放射工作人员应进行放射防护培训、个人剂量监测及职业健康检查管理。
\item 应及时向相关卫生健康行政部门申请放射性职业病危害预评价审核，取得卫生健康行政部门预评价审核同意批复后方可施工；建设项目在竣工验收前委托具有相应资质的放射卫生技术服务机构进行建设项目放射性职业病危害控制效果评价，完成竣工验收后及时办理《放射诊疗许可证》。
\item 如建设项目放射工作场所或辐射源项发生重大变更时应重新委托具有相应资质的放射卫生技术服务机构进行建设项目放射性职业病危害预评价。
\end{enumerate}
